% GENERATED FILE. Edit canonical lessons, then run npm run generate:books.
% canonical-source-hash: fnv1a64:cd3584108a205e5b
% canonical-lessons: UR-C259-pakvan, UR-C259-menyu, UR-C259-chatpata, UR-C259-maida, UR-C259-suji

\chapter{Dish, Menu, Tangy, Fine white flour, Semolina}
\label{ch:ur-dish-menu-tangy-fine-white-flour-semolina}
\hlchaptermodality{259}

\begin{chapteropening}
\textbf{By the end of this chapter:} \emph{I can say a dish, a menu, tangy, fine white flour and semolina.}

Complete the last lesson of chapter 259: I can say a dish, a menu, tangy, fine white flour and semolina.
\end{chapteropening}

\section[\texorpdfstring{\ur{پکوان} (pakvān)}{پکوان (pakvān)}]{\ur{پکوان} (pakvān) — a dish (of food)}

\label{lesson:UR-C259-pakvan}



\begin{quote}
Before the new one: say the Urdu for a tap, then the Urdu for a shower.
\end{quote}

\subsection*{You'll want to know: \ur{پکوان}}
\textbf{\ur{پکوان}} — \emph{pakvān} — \textquotedblleft{}a dish (of food)\textquotedblright{}.

\begin{cousinweb}
\textbf{\ur{پکوان}} (\emph{pakvān}) is usually taken from Sanskrit \emph{pakvānna}, \textquotedblleft{}cooked food\textquotedblright{}, built from \emph{pakva}, \textquotedblleft{}cooked\textquotedblright{}, and \emph{anna}, \textquotedblleft{}food\textquotedblright{}. CDIAL has no entry for it.
\end{cousinweb}

\begin{grammarlens}[title={What you've built}]
One more word for this chapter.
\end{grammarlens}

\subsection*{Your turn}
\begin{itemize}
\raggedright
  \item \emph{Say it:} \emph{pakvān}
  \item \emph{Say it:} \emph{pakvān}, once more
  \item \emph{Say it:} say \emph{pakvān}
  \item \emph{Recall:} say the Urdu for a tap, then the Urdu for a shower, then say \emph{pakvān} again
\end{itemize}

\begin{tcolorbox}[breakable,colback=teal!4,colframe=teal!35!black,title={Before you move on}]
What does \ur{پکوان} mean? (\textquotedblleft{}A dish (of food)\textquotedblright{}.) Say it once more.
\end{tcolorbox}
\section[\texorpdfstring{\ur{مینیو} (menyū)}{مینیو (menyū)}]{\ur{مینیو} (menyū) — a menu}

\label{lesson:UR-C259-menyu}



\begin{quote}
Before the new one: say the Urdu for a shower, then the Urdu for a dish.
\end{quote}

\subsection*{You'll want to know: \ur{مینیو}}
\textbf{\ur{مینیو}} — \emph{menyū} — \textquotedblleft{}a menu\textquotedblright{}.

\begin{cousinweb}
\textbf{\ur{مینیو}} (\emph{menyū}) is a loan from English \emph{menu}.
\end{cousinweb}

\begin{grammarlens}[title={What you've built}]
One more word for this chapter.
\end{grammarlens}

\subsection*{Your turn}
\begin{itemize}
\raggedright
  \item \emph{Say it:} \emph{menyū}
  \item \emph{Say it:} \emph{menyū}, once more
  \item \emph{Say it:} say \emph{menyū}
  \item \emph{Recall:} say the Urdu for a shower, then the Urdu for a dish, then say \emph{menyū} again
\end{itemize}

\begin{tcolorbox}[breakable,colback=teal!4,colframe=teal!35!black,title={Before you move on}]
What does \ur{مینیو} mean? (\textquotedblleft{}A menu\textquotedblright{}.) Say it once more.
\end{tcolorbox}
\section[\texorpdfstring{\ur{چٹپٹا} (chaṭpaṭā)}{چٹپٹا (chaṭpaṭā)}]{\ur{چٹپٹا} (chaṭpaṭā) — tangy}

\label{lesson:UR-C259-chatpata}



\begin{quote}
Before the new one: say the Urdu for a dish, then the Urdu for a menu.
\end{quote}

\subsection*{You'll want to know: \ur{چٹپٹا}}
\textbf{\ur{چٹپٹا}} — \emph{chaṭpaṭā} — \textquotedblleft{}tangy\textquotedblright{}.

\begin{cousinweb}
The origin of \textbf{\ur{چٹپٹا}} (\emph{chaṭpaṭā}) is not settled. CDIAL lists only a Hindi \emph{chaṭpaṭī}, \textquotedblleft{}restlessness\textquotedblright{}, under a reconstructed \emph{*chaṭṭa}, \textquotedblleft{}sudden movement\textquotedblright{} (CDIAL 4969), which may or may not be the same word.
\end{cousinweb}

\begin{grammarlens}[title={What you've built}]
One more word for this chapter.
\end{grammarlens}

\subsection*{Your turn}
\begin{itemize}
\raggedright
  \item \emph{Say it:} \emph{chaṭpaṭā}
  \item \emph{Say it:} \emph{chaṭpaṭā}, once more
  \item \emph{Say it:} say \emph{chaṭpaṭā}
  \item \emph{Recall:} say the Urdu for a dish, then the Urdu for a menu, then say \emph{chaṭpaṭā} again
\end{itemize}

\begin{tcolorbox}[breakable,colback=teal!4,colframe=teal!35!black,title={Before you move on}]
What does \ur{چٹپٹا} mean? (\textquotedblleft{}Tangy\textquotedblright{}.) Say it once more.
\end{tcolorbox}
\section[\texorpdfstring{\ur{میدہ} (maidā)}{میدہ (maidā)}]{\ur{میدہ} (maidā) — fine white flour}

\label{lesson:UR-C259-maida}



\begin{quote}
Before the new one: say the Urdu for a menu, then the Urdu for tangy.
\end{quote}

\subsection*{You'll want to know: \ur{میدہ}}
\textbf{\ur{میدہ}} — \emph{maidā} — \textquotedblleft{}fine white flour\textquotedblright{}.

\begin{cousinweb}
\textbf{\ur{میدہ}} (\emph{maidā}) is a loan, usually traced to Persian \emph{maida}, \textquotedblleft{}fine flour\textquotedblright{}.
\end{cousinweb}

\begin{grammarlens}[title={What you've built}]
One more word for this chapter.
\end{grammarlens}

\subsection*{Your turn}
\begin{itemize}
\raggedright
  \item \emph{Say it:} \emph{maidā}
  \item \emph{Say it:} \emph{maidā}, once more
  \item \emph{Say it:} say \emph{maidā}
  \item \emph{Recall:} say the Urdu for a menu, then the Urdu for tangy, then say \emph{maidā} again
\end{itemize}

\begin{tcolorbox}[breakable,colback=teal!4,colframe=teal!35!black,title={Before you move on}]
What does \ur{میدہ} mean? (\textquotedblleft{}Fine white flour\textquotedblright{}.) Say it once more.
\end{tcolorbox}
\section[\texorpdfstring{\ur{سوجی} (sūjī)}{سوجی (sūjī)}]{\ur{سوجی} (sūjī) — semolina}

\label{lesson:UR-C259-suji}



\begin{quote}
Before the new one: say the Urdu for tangy, then the Urdu for fine white flour.
\end{quote}

\subsection*{You'll want to know: \ur{سوجی}}
\textbf{\ur{سوجی}} — \emph{sūjī} — \textquotedblleft{}semolina\textquotedblright{}.

\begin{cousinweb}
\textbf{\ur{سوجی}} (\emph{sūjī}) is a tadbhava, a word inherited and reshaped by sound change. It comes from \emph{*sūjjī}, \textquotedblleft{}coarse wheat meal\textquotedblright{}, a form CDIAL reconstructs from its descendants (CDIAL 13552).
\end{cousinweb}

\begin{grammarlens}[title={What you've built}]
One more word for this chapter.
\end{grammarlens}

\subsection*{Your turn}
\begin{itemize}
\raggedright
  \item \emph{Say it:} \emph{sūjī}
  \item \emph{Say it:} \emph{sūjī}, once more
  \item \emph{Say it:} say \emph{sūjī}
  \item \emph{Recall:} say the Urdu for tangy, then the Urdu for fine white flour, then say \emph{sūjī} again
\end{itemize}

\begin{tcolorbox}[breakable,colback=teal!4,colframe=teal!35!black,title={Before you move on}]
What does \ur{سوجی} mean? (\textquotedblleft{}Semolina\textquotedblright{}.) Say it once more.
\end{tcolorbox}
